\section{Aufbau der Arbeit}

Die vorliegende Arbeit gliedert sich wie folgt: Zunächst wird in \textbf{Kapitel \ref{chapter:systemarchitektur}}  die \textbf{Systemarchitektur} des ReSiLENT-Systems dargelegt. Der Fokus liegt dabei auf der dezentralen Kommunikation via IPFS sowie den vier Detektionseinheiten, welche physikalische, systemnahe und netzwerkbasierte Anomalien auf Wallbox-Ebene erfassen.

\textbf{Kapitel \ref{chapter:thread_report}} beschreibt die \textbf{Konzeption des Threat Reporters}. Es wird erläutert, wie Detektionsergebnisse unter Anwendung zeitlicher Diskretisierung und Sliding-Window-Analysen in einem standardisierten CSV-Format aggregiert und im Peer-to-Peer-Netzwerk distribuiert werden.

In \textbf{Kapitel \ref{chapter:theoretische_grundlagen_der_informationsfusion}} werden die \textbf{theoretischen Grundlagen der Informationsfusion} adressiert. Angesichts variabler Eingabemengen wird die methodische Wahl einer Kombination aus \textit{Deep Sets} zur Gewährleistung der Permutationsinvarianz und \textit{Attention}-Mechanismen zur dynamischen Relevanzbewertung begründet.

Darauf aufbauend stellt \textbf{Kapitel \ref{chapter:methodik_gewichtete_agg_von_threat_reports}} die \textbf{Methodik der gewichteten Aggregation} vor. Hierbei werden die \textit{Negative Exponential Utility Function} sowie eine \textit{logistische Sigmoid-Funktion} hinsichtlich ihres asymptotischen Sättigungsverhaltens und ihrer Robustheit gegenüber statistischen Ausreißern evaluiert.

Die \textbf{Implementierung} und Integration dieser Modelle in den ReSiLENT-Datenfluss werden in \textbf{Kapitel \ref{chapter:evaluation}} beschrieben. Die resultierenden aggregierten CTI-Reports werden präsentiert und hinsichtlich ihrer Qualität für die knoteninterne Nutzung bewertet.

Abschließend erfolgt in \textbf{Kapitel \ref{chapter:diskussion_und_forschungsluecke}} eine \textbf{Diskussion der Ergebnisse}. Dabei werden kritische Forschungslücken identifiziert, insbesondere die empirische Validierung der Gewichtungsstrategien, das Potenzial LLM-gestützter Aufbereitung für SOC-Operatoren sowie die Untersuchung systemischer Risiken wie des \textit{Echo-Chamber-Effekts}.
